%% Question bank chapter file - generated by bank_tools.py build (do not edit while the bank is still being built:
%% it is regenerated from the per-test data in Question Bank\data). Every question is one bankq block:
%%   \begin{bankq}[id=..., source={...}, paper=.., qno=.., type=..., subtopic={...}, answer={...}, answer_src=key|solved,
%%                 working={...}, difficulty=1-10, diff_parts={K S I C F}, diff_note={...},
%%                 group=<passage id or empty>, also_in={duplicate copies in other tests}]
%% difficulty = min(10, K+S+I+C+F-1): K concepts 1-3, S steps 1-4, I insight 0-2, C computation 0-2, F format 0-1
%%   ...question LaTeX, options, figures (\bfig / \bfigi -> img/)...
%%   \end{bankq}
%% Passages shared by several questions are bankpassage blocks; their questions carry group=<passage id>.
\documentclass[10pt]{article}
\usepackage[a4paper,margin=1.8cm]{geometry}
\usepackage[T1]{fontenc}
\usepackage{amsmath,amssymb,mathtools}
\usepackage{graphicx}
\usepackage[export]{adjustbox}
\usepackage{xcolor,needspace,keyval}
\usepackage{enumitem}
\usepackage{booktabs,array,tabularx}
\usepackage[hidelinks]{hyperref}
\setlength{\parindent}{0pt}
\setlength{\parskip}{3pt}
\makeatletter
\define@key{bq}{id}{\def\bq@id{#1}}\define@key{bq}{source}{\def\bq@src{#1}}\define@key{bq}{paper}{\def\bq@paper{#1}}
\define@key{bq}{qno}{\def\bq@qno{#1}}\define@key{bq}{type}{\def\bq@type{#1}}\define@key{bq}{subtopic}{\def\bq@sub{#1}}
\define@key{bq}{answer}{\def\bq@ans{#1}}\define@key{bq}{answer_src}{\def\bq@asrc{#1}}\define@key{bq}{working}{\def\bq@work{#1}}
\define@key{bq}{group}{\def\bq@group{#1}}\define@key{bq}{also_in}{\def\bq@also{#1}}
\define@key{bq}{difficulty}{\def\bq@diff{#1}}\define@key{bq}{diff_parts}{\def\bq@dparts{#1}}\define@key{bq}{diff_note}{\def\bq@dnote{#1}}
\newenvironment{bankq}[1][]{\def\bq@work{}\def\bq@also{}\def\bq@group{}\def\bq@diff{?}\def\bq@dparts{}\def\bq@dnote{}\setkeys{bq}{#1}%
  \par\needspace{4\baselineskip}\medskip\noindent{\small\textbf{\bq@id}\hfill\color{gray}[\bq@src, Paper \bq@paper, Q\bq@qno]
  \bq@type\ $\cdot$ \bq@sub\ $\cdot$ \textbf{D\bq@diff}}\par\nopagebreak\relax}
  {\par{\small\color{teal}Answer: \textbf{\bq@ans}\ifx\bq@work\@empty\else\ (solved: \bq@work)\fi
  \ifx\bq@also\@empty\else\par\color{gray}Also in: \bq@also\fi}\par}
\define@key{bp}{id}{\def\bp@id{#1}}\define@key{bp}{source}{\def\bp@src{#1}}\define@key{bp}{questions}{\def\bp@qs{#1}}
\newenvironment{bankpassage}[1][]{\setkeys{bp}{#1}\par\needspace{4\baselineskip}\bigskip\noindent
  {\small\textbf{Passage \bp@id}\hfill\color{gray}[\bp@src; questions \bp@qs]}\par\nopagebreak\relax\itshape}{\par}
\makeatother
\newcommand{\bfig}[2]{\par\smallskip{\centering\includegraphics[width=#1pt,max width=0.62\linewidth,max height=0.32\textheight]{img/#2}\par}\smallskip}
\newcommand{\bfigi}[2]{\includegraphics[width=#1pt,max width=\linewidth,max height=0.32\textheight,valign=c]{img/#2}}
\newcommand{\bmissing}[1]{\textit{[figure missing in the source conversion: \detokenize{#1}]}}
\newenvironment{qoptions}{\begin{itemize}[leftmargin=1.8em,itemsep=1pt,topsep=2pt,label={}]}{\end{itemize}}
\newcommand{\opt}[2]{\item[(#1)] #2}
\providecommand{\fwstart}{}\providecommand{\fwend}{}\providecommand{\qfullsplit}{}
\newcolumntype{L}[1]{>{\raggedright\arraybackslash}p{#1}}
\newcolumntype{C}[1]{>{\centering\arraybackslash}p{#1}}
\newcolumntype{R}[1]{>{\raggedleft\arraybackslash}p{#1}}
\providecommand{\sectionhdr}[1]{}\providecommand{\subsectionhdr}[1]{}


\begin{document}
{\Large\bfseries WPE --- Question Bank}\par\smallskip
{\small Physics chapter P03. Questions from all converted test papers, grouped by source series. Answers: printed key where it matches the paper, otherwise solved (answer\_src=solved, with working). Exact repeats across papers appear once, with also\_in listing the other copies.}\par\medskip\hrule

\section*{Score-I papers}
\begin{bankq}[id=SI12-P1-PHY-S2-Q05, source={Score-I - Test 12 (09 Mar 2022)}, paper=1, qno=5, type=numerical, subtopic={Work}, answer={2.00}, answer_src=key, difficulty=2, diff_parts={K1 S1 I0 C0 F1}, diff_note={concepts: work-energy theorem | steps: W = F.d; final kinetic energy gives k}]
A particle of mass 2 kg moves on a smooth horizontal plane under the action of a single force $\vec F=(3\hat i+4\hat j)$ N. Under this force it is displaced from $(0,0)$ to $(1\text{ m},1\text{ m})$. The initial speed of the particle is $\sqrt3$ m/s and final speed is $\sqrt{5k}$ m/s. Find the value of $k$.
\end{bankq}
\begin{bankpassage}[id=SI12-P2-PHY-S1-G1-2, source={Score-I - Test 12 (09 Mar 2022)}, questions={1-2}]
In column-1, the variation of potential energy (in J) as a function of $x$ is given for a particle. The particle moves under the effect of conservative force acting along the $x$-axis. Assume that no other force is acting on the particle. In column-2, locations on the $x$-axis are mentioned, and in column-3, various choices for the kinetic energy of particle are given.
\begin{center}\footnotesize
\begin{tabular}{|p{0.7cm}|p{5.2cm}|p{0.7cm}|p{2.2cm}|p{0.7cm}|p{2.4cm}|}
\hline
 & \centering Column--1\newline Potential energy as function of $x$ & \multicolumn{2}{c|}{\shortstack{Column--2\\Location on the $x$-axis}} & \multicolumn{2}{c|}{\shortstack{Column--3\\Kinetic energy of particle (in J)}}\\
\hline
(I) & $U(x)=\dfrac1{100}\left[x^3-6x^2+11x-6\right]$ & (i) & $x=0$ & (P) & KE $=0$\\
(II) & $U(x)=-2x^2+8x-6$ & (ii) & $x=1$ & (Q) & KE $=1$\\
(III) & $U(x)=\dfrac13\left[x^4-7x^3+17x^2-17x+6\right]$ & (iii) & $x=2$ & (R) & KE $=2$\\
(IV) & $U(x)=2x^2-4x$ & (iv) & $x=3$ & (S) & KE $=3$\\
\hline
\end{tabular}
\end{center}
\end{bankpassage}
\begin{bankq}[id=SI12-P2-PHY-S1-Q01, source={Score-I - Test 12 (09 Mar 2022)}, paper=2, qno=1, type=single, subtopic={Potential energy curves}, answer={C}, answer_src=key, difficulty=3, diff_parts={K2 S2 I0 C0 F0}, diff_note={concepts: force from a potential energy curve, mechanical energy | steps: read slope sign in each region; combine with allowed region for E = 3 J; match the three lists}, group=SI12-P2-PHY-S1-G1-2]
In which of the following case the force directed towards $+$ve $x$-axis if mechanical energy of the particle is 3 J?
\begin{qoptions}
\opt{A}{(II) (ii) (S)}
\opt{B}{(I) (iv) (S)}
\opt{C}{(II) (iv) (S)}
\opt{D}{(II) (iii) (Q)}
\end{qoptions}
\end{bankq}
\begin{bankq}[id=SI12-P2-PHY-S1-Q02, source={Score-I - Test 12 (09 Mar 2022)}, paper=2, qno=2, type=single, subtopic={Potential energy curves}, answer={A}, answer_src=key, difficulty=3, diff_parts={K2 S2 I0 C0 F0}, diff_note={concepts: potential energy curve, energy conservation | steps: locate minima for each curve; use E = 3 J to check the allowed region; match}, group=SI12-P2-PHY-S1-G1-2]
In which of the following the particle is in stable equilibrium if mechanical energy of the particle is 3 J?
\begin{qoptions}
\opt{A}{(III) (ii) (S)}
\opt{B}{(II) (iii) (Q)}
\opt{C}{(III) (i) (Q)}
\opt{D}{(II) (ii) (S)}
\end{qoptions}

\noindent
\end{bankq}

\section*{TOAS papers}
\begin{bankq}[id=TOAS04-P1-PHY-S1-Q03, source={TOAS - Test 04 (15 Mar 2022)}, paper=1, qno=3, type=single, subtopic={Energy}, answer={C}, answer_src=key, difficulty=3, diff_parts={K2 S2 I0 C0 F0}, diff_note={concepts: spring force on a rubber thread, energy | steps: k from the lower ball just leaving the ground; work done by the agent as spring energy plus gravity change; speed after release}]
Two tennis balls of mass 60 g are attached with a massless rubber thread, and held in the vertical position as shown in the figure. In this position, the unstretched length of the rubber thread is 40 cm. The upper ball is slowly raised vertically upward, until the lower ball just becomes unsupported by the ground. At this time the length of the thread is 1 m. The rubber thread exerts a force which is proportional to its extension. Choose correct statement.
\bfig{50.6}{TOAS04_p1_phy_s1_q03_1.png}
\begin{qoptions}
\opt{A}{Work done by external agent while the upper ball was raised is 0.6 J.}
\opt{B}{Work done by external agent while the upper ball was raised is 0.23 J.}
\opt{C}{Releasing the upper ball, the speed which it hits the lower one is 5.1 m/s.}
\opt{D}{The time that elapses between the release of the upper ball and the collision is 0.5 sec.}
\end{qoptions}
\end{bankq}
\begin{bankq}[id=TOAS04-P1-PHY-S1-Q04, source={TOAS - Test 04 (15 Mar 2022)}, paper=1, qno=4, type=single, subtopic={Power}, answer={B}, answer_src=key, difficulty=3, diff_parts={K1 S2 I1 C0 F0}, diff_note={concepts: constant power motion | steps: a = P/(m v) so a halves when v doubles; work-energy P t = delta KE; t = 9/4 s | I: acceleration halves when speed doubles}]
A mass of $m=1$ kg moves initially with a velocity of $v_0=10\sqrt{15}$ m/s. A constant power $P=1000$ Watt is applied so as to increase its velocity. Find the time that elapses before the acceleration is reduced to half of its original value.
\begin{qoptions}
\opt{A}{4 sec}
\opt{B}{9/4 sec}
\opt{C}{1 sec}
\opt{D}{4.5 sec}
\end{qoptions}
\end{bankq}
\begin{bankq}[id=TOAS04-P1-PHY-S2-Q06, source={TOAS - Test 04 (15 Mar 2022)}, paper=1, qno=6, type=multi, subtopic={Energy}, answer={D}, answer_src=key, difficulty=2, diff_parts={K1 S1 I0 C0 F1}, diff_note={concepts: projectile from a sliding bead | steps: speed at the bottom v\^{}2 = 2 g R; D\^{}2 = v\^{}2 (2H/g) gives RH proportional to D\^{}2}]
A semicircular wire of radius $R$ is oriented vertically. A small bead is released from rest from the top of the wire. It slides without friction under the influence of gravity to the bottom, where it then leaves the wire horizontally and falls distance $H$ to the ground. The bead lands a horizontal distance $D$ away from where it was launched. Which of the following is correct graph of $RH$ versus $D^2$ ?
\bfig{74.7}{TOAS04_p1_phy_s2_q06_1.png}
\begin{qoptions}
\opt{A}{\bfigi{80.6}{TOAS04_p1_phy_s2_q06_A.png}}
\opt{B}{\bfigi{87.0}{TOAS04_p1_phy_s2_q06_B.png}}
\opt{C}{\bfigi{84.5}{TOAS04_p1_phy_s2_q06_C.png}}
\opt{D}{\bfigi{85.0}{TOAS04_p1_phy_s2_q06_D.png}}
\end{qoptions}
\end{bankq}
\begin{bankq}[id=TOAS04-P2-PHY-S1-Q03, source={TOAS - Test 04 (15 Mar 2022)}, paper=2, qno=3, type=multi, subtopic={Work}, answer={A,C}, answer_src=key, difficulty=3, diff_parts={K1 S1 I1 C0 F1}, diff_note={concepts: work done by friction | steps: friction on the man is forward but his displacement relative to the ground; friction on the plank | I: work depends on displacement of each body}]
A man is running with a constant acceleration (without slipping) on a plank which is placed on a horizontal smooth surface. Then choose the correct option :-
\begin{qoptions}
\opt{A}{Work done by friction on the man is negative.}
\opt{B}{Work done by friction on the man is positive.}
\opt{C}{Work done by friction on the plank is positive.}
\opt{D}{Work done by friction on the plank is negative.}
\end{qoptions}
\end{bankq}
\begin{bankq}[id=TOAS05-P1-PHY-S1-Q04, source={TOAS - Test 05 (03 Apr 2022)}, paper=1, qno=4, type=multi, subtopic={Energy}, answer={B,D}, answer_src=key, difficulty=5, diff_parts={K2 S2 I1 C0 F1}, diff_note={concepts: energy conservation, projectile from a hill | steps: v\^{}2 = 2g(H - y); fall time; x as a function of y and its maximum | I: x = 2 sqrt(y(H - y))}]
A small block slides down a hill (starting from rest) of height $H$. The part of the hill near its end $B$ is horizontal. The block flies off the hill at $B$ to hit the ground at $C$. Hill is smooth. If $y$ is height of end $B$ of hill from ground and $x$ is horizontal displacement of block before hitting the ground after loosing contact with track, then :
\bfig{152.0}{TOAS05_p1_phy_s1_q04_1.png}
\begin{qoptions}
\opt{A}{For $x$ to be maximum, $y$ is $H$}
\opt{B}{For $x$ to be maximum $y$ is $\dfrac H2$.}
\opt{C}{Maximum possible value of $x$ is $\dfrac H2$.}
\opt{D}{Maximum possible value of $x$ is $H$.}
\end{qoptions}
\end{bankq}
\begin{bankq}[id=TOAS05-P1-PHY-S3-Q01, source={TOAS - Test 05 (03 Apr 2022)}, paper=1, qno=1, type=numerical, subtopic={Work}, answer={18.00}, answer_src=key, difficulty=2, diff_parts={K1 S1 I0 C0 F1}, diff_note={concepts: work of friction on a sliding chain | steps: normal reaction of the resting portion as a function of y; integrate friction from 2L/3 to L}]
A uniform chain of mass $M$ and length $L$ rests on a rough horizontal table such that one of its ends hangs over the edge of the table. Chain begins to slide if length of overhanging part is at least equal to two third of the total length of the chain. Chain is released and starts its motion from a position where two third of its total length is overhanging. If work done by friction force acting on the chain by the moment it completely slides off the table is given by $\dfrac{-\mu MgL}{N}$ then find the value of $N$. Coefficient of friction between chain and rough horizontal surface is $\mu$ and acceleration due to gravity is $g$.
\bfig{78.1}{TOAS05_p1_phy_s3_q01_1.png}
\end{bankq}
\begin{bankq}[id=TOAS05-P2-PHY-S1-Q07, source={TOAS - Test 05 (03 Apr 2022)}, paper=2, qno=7, type=single, subtopic={Power}, answer={A}, answer_src=key, difficulty=1, diff_parts={K1 S1 I0 C0 F0}, diff_note={concepts: power of a pump | steps: power used in lifting; the remainder gives kinetic energy per second}]
A pump of 200 W is lifting 2kg water per second from an average depth of 10m. The velocity with which water comes out of pump is :- (Take : $g=9.8$ m/s$^2$)
\begin{qoptions}
\opt{A}{2 m/s}
\opt{B}{2.5 m/s}
\opt{C}{4 m/s}
\opt{D}{3 m/s}
\end{qoptions}
\end{bankq}
\begin{bankq}[id=TOAS05-P2-PHY-S1-Q08, source={TOAS - Test 05 (03 Apr 2022)}, paper=2, qno=8, type=single, subtopic={Energy}, answer={B}, answer_src=key, difficulty=2, diff_parts={K1 S1 I1 C0 F0}, diff_note={concepts: energy of a system with springs | steps: work done by gravity; energy stored in the springs of the arrangement | I: effective spring constant}]
A system is released from rest as shown in figure. Kinetic energy of mass $m$ when its moves distance $x$ is downward direction is. (initially both spring are unstretched and all spring are massless)
\bfig{33.7}{TOAS05_p2_phy_s1_q08_1.png}
\begin{qoptions}
\opt{A}{$10\,mgx$}
\opt{B}{$\dfrac{10mgx-kx^2}{10}$}
\opt{C}{$mgx-kx^2$}
\opt{D}{$mgx-\dfrac{kx^2}2$}
\end{qoptions}
\end{bankq}
\begin{bankq}[id=TOAS05-P2-PHY-S1-Q10, source={TOAS - Test 05 (03 Apr 2022)}, paper=2, qno=10, type=single, subtopic={Work}, answer={A}, answer_src=key, difficulty=1, diff_parts={K1 S1 I0 C0 F0}, diff_note={concepts: work-displacement graph | steps: integrate the F-x graph piece by piece}]
A particle is displaced form $x=-6$m to $x=+6$m. A force $F$ acting on the particle during its motion is shown in figure. Graph between work done by this force ($W$) and displacement ($x$) should be :-
\bfig{84.2}{TOAS05_p2_phy_s1_q10_1.png}
\begin{qoptions}
\opt{A}{\bfigi{69.3}{TOAS05_p2_phy_s1_q10_A.png}}
\opt{B}{\bfigi{74.0}{TOAS05_p2_phy_s1_q10_B.png}}
\opt{C}{\bfigi{79.9}{TOAS05_p2_phy_s1_q10_C.png}}
\opt{D}{\bfigi{75.2}{TOAS05_p2_phy_s1_q10_D.png}}
\end{qoptions}
\end{bankq}
\begin{bankq}[id=TOAS05-P2-PHY-S1-Q12, source={TOAS - Test 05 (03 Apr 2022)}, paper=2, qno=12, type=single, subtopic={Work}, answer={B}, answer_src=key, difficulty=1, diff_parts={K1 S1 I0 C0 F0}, diff_note={concepts: work by a position-dependent force | steps: parametrise the line; integrate F dot dr}]
A force acting on a particle moving in the $xy$-plane is given by $\vec F=\left(2y\hat i+x^2\hat j\right)$ N, where $x$ and $y$ are in m. The particle moves along a straight line from the origin to $(5,5)$. The work done by $F$ is :-
\begin{qoptions}
\opt{A}{125 J}
\opt{B}{66.7 J}
\opt{C}{35 J}
\opt{D}{25 J}
\end{qoptions}
\end{bankq}
\begin{bankq}[id=TOAS05-P2-PHY-S1-Q14, source={TOAS - Test 05 (03 Apr 2022)}, paper=2, qno=14, type=single, subtopic={Conservative forces}, answer={B}, answer_src=key, difficulty=1, diff_parts={K1 S1 I0 C0 F0}, diff_note={concepts: force from potential energy | steps: F = -grad U}]
A particle in a certain conservative force field has a potential energy given by $U=\dfrac{20yz}{x}$. The force exerted on it is :-
\begin{qoptions}
\opt{A}{$\left(\dfrac{20yz}{x^2}\right)\hat i+\left(\dfrac{20z}x\right)\hat j+\left(\dfrac{20y}x\right)\hat k$}
\opt{B}{$\left(\dfrac{20yz}{x^2}\right)\hat i-\left(\dfrac{20z}x\right)\hat j-\left(\dfrac{20y}x\right)\hat k$}
\opt{C}{$-\left(\dfrac{20yz}{x^2}\right)\hat i-\left(\dfrac{20z}x\right)\hat j-\left(\dfrac{20y}x\right)\hat k$}
\opt{D}{$-\left(\dfrac{20yz}{x^2}\right)\hat i+\left(\dfrac{20z}x\right)\hat j+\left(\dfrac{20y}x\right)\hat k$}
\end{qoptions}
\end{bankq}
\begin{bankq}[id=TOAS05-P2-PHY-S2-Q09, source={TOAS - Test 05 (03 Apr 2022)}, paper=2, qno=9, type=numerical, subtopic={Work}, answer={10}, answer_src=key, difficulty=2, diff_parts={K1 S1 I0 C0 F1}, diff_note={concepts: work along a path | steps: work on OA and work on AB; add}]
A variable force $\vec F=\left(x^2\hat i+xy\hat j\right)$ N is acting on a particle which moves from point $(0,0)$ to $(1\mathrm m,2\mathrm m)$ along the path $OAB$ as shown. The work done by force is $\dfrac pq$ joules, where $p$ and $q$ are co-prime integers. Find the value of $(p+q)$.
\bfig{120.8}{TOAS05_p2_phy_s2_q09_1.png}
\end{bankq}
\begin{bankq}[id=TOAS06-P1-PHY-S3-Q03, source={TOAS - Test 06 (17 Apr 2022)}, paper=1, qno=3, type=numerical, subtopic={Work}, answer={15.00}, answer_src=key, difficulty=2, diff_parts={K1 S1 I0 C0 F1}, diff_note={concepts: work-energy with a nonlinear spring | steps: spring work as an integral of k x\^{}2; equate to kinetic energy}]
Design considerations for the bumper B on the 5000 kg train car require use of a nonlinear spring having the load-deflection characteristics shown in the graph. Select the proper value of $k$ (in MN/m$^2$) so that the maximum deflection of the spring is limited to 0.2 m when the car, traveling at 4 m/s, strikes the rigid stop. Neglect the mass of the car wheels.
\bfig{273.8}{TOAS06_p1_phy_s3_q03_1.png}
\end{bankq}
\begin{bankq}[id=TOAS13-P1-PHY-S1-Q02, source={TOAS - Test 13 (24 Jul 2022)}, paper=1, qno=2, type=multi, subtopic={Work}, answer={C}, answer_src=key, difficulty=5, diff_parts={K2 S2 I1 C0 F1}, diff_note={concepts: work by a variable force, work-energy theorem | steps: force on the rod grows as more length enters the field; work over the first l is mgl/2; remaining 4l at full weight; v from work-energy}]
Consider a situation in which a uniform gravitation field is present below line AB and no gravity above line AB. Then the speed with which a rod shown in the figure will strike earth surface when released from the position shown is
\bfig{100.4}{TOAS13_p1_phy_s1_q02_1.png}
\begin{qoptions}
\opt{A}{$\sqrt{g\ell}$}
\opt{B}{$2\sqrt{g\ell}$}
\opt{C}{$3\sqrt{g\ell}$}
\opt{D}{$4\sqrt{g\ell}$}
\end{qoptions}
\end{bankq}
\begin{bankq}[id=TOAS17-P2-PHY-S2-Q03, source={TOAS - Test 17 (18 Sep 2022)}, paper=2, qno=3, type=numerical, subtopic={Potential energy curves}, answer={4}, answer_src=key, difficulty=6, diff_parts={K2 S2 I1 C1 F1}, diff_note={concepts: equilibrium position, small oscillation about a minimum | steps: equilibrium from dU/dx; second derivative; period; compare with the form}]
The potential energy of a particle of mass $m$ free to move along $x$ axis has the form $U(x)=\dfrac a{x^2}-\dfrac bx$, where $a$ and $b$ are positive constant and $x$ is $x$-coordinate of the particle. The time period of SHM (along $x$-axis) of small amplitude around mean position is $\dfrac{P\pi a}{b^2}\sqrt{2ma}$. Find the value of $P$.
\end{bankq}
\begin{bankq}[id=TOAS19-P1-PHY-S1-Q02, source={TOAS - Test 19 (05 Feb 2023)}, paper=1, qno=2, type=multi, subtopic={Potential energy curves}, answer={A,D}, answer_src=key, difficulty=4, diff_parts={K1 S2 I1 C0 F1}, diff_note={concepts: equilibrium and oscillation | steps: stable equilibrium need not give simple harmonic motion; unstable equilibrium motion; counterexamples}]
Choose the \textit{INCORRECT} statement(s) :-
\begin{qoptions}
\opt{A}{All the system in stable equilibrium, if disturbed perform simple harmonic motion.}
\opt{B}{All the system in unstable equilibrium, if disturbed do not perform simple harmonic motion.}
\opt{C}{All the system in stable equilibrium, if disturbed perform oscillatory motion.}
\opt{D}{Any system in unstable equilibrium, if disturbed will not perform oscillatory motion.}
\end{qoptions}
\end{bankq}
\begin{bankq}[id=TOAS19-P1-PHY-S3-Q01, source={TOAS - Test 19 (05 Feb 2023)}, paper=1, qno=1, type=numerical, subtopic={Power}, answer={2.25}, answer_src=key, difficulty=3, diff_parts={K1 S1 I1 C0 F1}, diff_note={concepts: constant power motion | steps: a = P/mv so a halves when v doubles; work-energy gives the time}]
A mass 1kg moves initially with a velocity $v_0=10\sqrt{15}$ m/s. A constant power $P=1000$ W is applied so as to increase velocity find the time in seconds after which acceleration is reduced to half of its original value.
\bfig{82.0}{TOAS19_p1_phy_s3_q01_1.png}
\end{bankq}

\section*{TOPS papers}
\begin{bankq}[id=TOPS02-P1-PHY-S1-Q02, source={TOPS - Test 02 (28 Aug 2022)}, paper=1, qno=2, type=single, subtopic={Work}, answer={D}, answer_src=key, difficulty=2, diff_parts={K1 S1 I0 C1 F0}, diff_note={concepts: work done by forces | steps: work by static friction and normal; check each case}]
In which of the following cases is the work done non-zero?
\begin{qoptions}
\opt{A}{Work done by normal of ladder on you when you climb a ladder}
\opt{B}{Work done by friction on your shoes when you run in a race}
\opt{C}{Work done by a pencil on a paper when you write on the paper}
\opt{D}{None of the above}
\end{qoptions}
\end{bankq}
\begin{bankq}[id=TOPS02-P1-PHY-S2-Q11, source={TOPS - Test 02 (28 Aug 2022)}, paper=1, qno=11, type=multi, subtopic={Work}, answer={B,C,D}, answer_src=key, difficulty=6, diff_parts={K2 S2 I1 C1 F1}, diff_note={concepts: work done by friction, kinematics | steps: friction grows as the sledge enters; average work half; deceleration after fully on; distances}]
A sledge moving over a smooth horizontal surface of ice at a velocity $v_0$ drives out on a horizontal road and comes to a halt as shown. The sledge has a length $\ell$, mass $m$ and friction between sledge and road is $\mu$.
\bfig{119.0}{TOPS02_p1_phy_s2_q11_1.png}
\begin{qoptions}
\opt{A}{No work is done by the friction to switch the sledge from ice to the road.}
\opt{B}{A work of $\dfrac12\mu mg\ell$ is done against friction while sledge switches completely on to road.}
\opt{C}{The distance covered by the sledge on the road is $\left(\dfrac{v_0^2}{2\mu g}-\dfrac\ell2\right)$.}
\opt{D}{Total distance moved by the sledge before stopping is $\left(\dfrac{v_0^2}{2\mu g}+\dfrac\ell2\right)$.}
\end{qoptions}
\end{bankq}
\begin{bankq}[id=TOPS03-P1-PHY-S1-Q01, source={TOPS - Test 03 (20 Nov 2022)}, paper=1, qno=1, type=single, subtopic={Work}, answer={B}, answer_src=key, difficulty=9, diff_parts={K3 S3 I2 C2 F0}, diff_note={concepts: friction, work-energy, integration | steps: friction at each finger as it moves; which finger slides; normal forces as a function of position; work by the external agent as an integral; sum over the stages | alternating sliding fingers}]
A uniform ruler of mass $m$ and length $\ell$ is balanced horizontally on the two index fingers as shown in the figure (top view). Initially the fingers are at the ends of the ruler. Now, one hand at a time, the index fingers are moved slowly in horizontal direction towards the other to meet at the middle of the ruler, during which the ruler slowly slides one way or the other. If the coefficients of static and kinetic frictions between the ruler and fingers are $\mu$ and $\mu/3$ respectively, then what is total work done by the external agent till the hands meet?
\bfig{91.8}{TOPS03_p1_phy_s1_q01_1.png}
\begin{qoptions}
\opt{A}{$\dfrac{\mu mg\ell}4$}
\opt{B}{$\dfrac{\mu mg\ell}6\ln\left(\dfrac{3\sqrt3}2\right)$}
\opt{C}{$\dfrac{\mu mg\ell}2\ln\left(\dfrac52\right)$}
\opt{D}{None of these}
\end{qoptions}
\end{bankq}
\begin{bankq}[id=TOPS03-P1-PHY-S1-Q04, source={TOPS - Test 03 (20 Nov 2022)}, paper=1, qno=4, type=single, subtopic={Power}, answer={C}, answer_src=key, difficulty=5, diff_parts={K2 S2 I1 C1 F0}, diff_note={concepts: relative velocity, friction power | steps: relative velocity vector of A and B; kinetic friction; power equals force times relative speed; value | vector relative velocity}]
A small block $A$ of mass 1 kg is placed over a horizontal plate $B$ of mass 2 kg which is placed on a smooth horizontal surface. Block $A$ is being moved continuously 1 m/s as shown in the figure and the plate is moved continuously with speed 1 m/s, towards the direction shown in the figure. If the coefficient of friction between blocks is 0.1, the net power loss due to friction is : ($g=10$ m/s$^2$)
\bfig{69.4}{TOPS03_p1_phy_s1_q04_1.png}
\begin{qoptions}
\opt{A}{0.5 J/s}
\opt{B}{1 J/s}
\opt{C}{$\sqrt2$ J/s}
\opt{D}{2 J/s}
\end{qoptions}
\end{bankq}
\begin{bankq}[id=TOPS03-P2-PHY-S1-Q05, source={TOPS - Test 03 (20 Nov 2022)}, paper=2, qno=5, type=multi, subtopic={Energy}, answer={A,B}, answer_src=key, difficulty=8, diff_parts={K3 S2 I2 C1 F1}, diff_note={concepts: energy conservation, spring, constraint | steps: equilibrium extension; energy conservation for the maximum extension; velocity at the equilibrium; check each | spring with two particles}]
Particles $A$ and $B$ are connected by a massless, inextensible cord which goes around the fixed massless pulley. If the system is released from rest with the spring initially unstretched, then:
\bfig{63.2}{TOPS03_p2_phy_s1_q05_1.png}
\begin{qoptions}
\opt{A}{The maximum extension in the spring is $\dfrac{2mg}K$.}
\opt{B}{The maximum velocity of particle $C$ is $\sqrt{\dfrac{3m}{4K}}g$.}
\opt{C}{The maximum extension in the spring is $\dfrac{mg}K$.}
\opt{D}{The maximum velocity of particle $C$ is $\sqrt{\dfrac{3m}{2K}}g$.}
\end{qoptions}
\end{bankq}
\begin{bankq}[id=TOPS03-P2-PHY-S2-Q06, source={TOPS - Test 03 (20 Nov 2022)}, paper=2, qno=6, type=numerical, subtopic={Work}, answer={0.50}, answer_src=key, difficulty=7, diff_parts={K2 S2 I2 C1 F1}, diff_note={concepts: work done by a non conservative force, area | steps: work equals the line integral; Green theorem; multiply by the area}]
In the $x$-$y$ plane there exist a closed figure of total area 25 cm$^2$. Find the magnitude of work done by the $\vec F=K\left(y\hat i-x\hat j\right)$ in one complete round around the circumference of the figure. Take $K=100$N/m.
\bfig{93.0}{TOPS03_p2_phy_s2_q06_1.png}
\end{bankq}
\begin{bankq}[id=TOPS04-P1-PHY-S1-Q05, source={TOPS - Test 04 (18 Dec 2022)}, paper=1, qno=5, type=multi, subtopic={Work}, answer={A,C,D}, answer_src=key, difficulty=6, diff_parts={K2 S2 I1 C1 F1}, diff_note={concepts: belt friction, work | steps: relative motion on the belt; kinetic friction phase; static phase; work and acceleration at 6 s}]
A block of mass 10 kg is put gently on a belt-conveyor system of infinite length at $t=0$ sec, which is moving with constant speed 20 m/sec rightward at all time, irrespectively of any situation by means of a motor-system as shown in the figure. Constant force of magnitude 15 N is applied on the block continuously during its motion. Mark the INCORRECT statement(s) :
\bfig{213.9}{TOPS04_p1_phy_s1_q05_1.png}
\begin{qoptions}
\opt{A}{The nature of friction force acting on block of mass 10 kg during its motion is always kinetic in nature.}
\opt{B}{Work done by the kinetic friction on the block of mass 10 kg is 1250 Joule.}
\opt{C}{The magnitude of acceleration of the block of mass 10 kg at $t=6$ sec is 4 m/s$^2$.}
\opt{D}{The magnitude of friction on the block of mass 10 kg at $t=6$ sec is zero.}
\end{qoptions}
\end{bankq}
\begin{bankpassage}[id=TOPS04-P1-PHY-S3-G3-4, source={TOPS - Test 04 (18 Dec 2022)}, questions={3-4}]
A small block of mass $m$ slides down an inclined plane forming an angle $\alpha=30^\circ$ with the horizontal. The coefficient of friction $k$ between the block and the plane is no longer constant but varies with $x$, the distance travelled by the block on the plane, according to the equation $k=ax$, where $a$ is a constant. (Use : $a=\dfrac{\sqrt3}2$, $g=10$ m/s$^2$)
\end{bankpassage}
\begin{bankq}[id=TOPS04-P1-PHY-S3-Q03, source={TOPS - Test 04 (18 Dec 2022)}, paper=1, qno=3, type=numerical, subtopic={Work}, answer={1.335}, answer_src=key, difficulty=7, diff_parts={K2 S2 I2 C1 F1}, diff_note={concepts: variable friction, work energy | steps: net force along the incline; work integral; stopping condition; distance}, group=TOPS04-P1-PHY-S3-G3-4]
Find the distance (in meter) travelled by the block before coming to rest.
\end{bankq}
\begin{bankq}[id=TOPS04-P1-PHY-S3-Q04, source={TOPS - Test 04 (18 Dec 2022)}, paper=1, qno=4, type=numerical, subtopic={Work}, answer={1.825}, answer_src=key, difficulty=7, diff_parts={K2 S2 I2 C1 F1}, diff_note={concepts: variable friction, work energy | steps: force zero at maximum speed; position; work to that point; speed}, group=TOPS04-P1-PHY-S3-G3-4]
Find the maximum speed (in m/s) of block during the motion.
\end{bankq}
\begin{bankpassage}[id=TOPS04-P1-PHY-S3-G5-6, source={TOPS - Test 04 (18 Dec 2022)}, questions={5-6}]
A weightless rod of length $\ell$ with a small load of mass $m$ at one of its end is held vertical with its lower end hinged on a horizontal surface. The load touches a wedge of mass $M$ in this position. A slight jerk towards right sets the system in motion (see figure), with rod rotating freely in vertical plane about its lower end. There is no friction. (Use : $g\ell=16$ SI units)
\bfig{157.0}{TOPS04_p1_phy_s3_para5_6_1.png}
\end{bankpassage}
\begin{bankq}[id=TOPS04-P1-PHY-S3-Q05, source={TOPS - Test 04 (18 Dec 2022)}, paper=1, qno=5, type=numerical, subtopic={Energy}, answer={4.00}, answer_src=solved, working={Load speed $v$, wedge speed $u=v\cos\theta$; energy gives $u^2=\dfrac{2g\ell c^2(1-c)}{1+\mu c^2}$ with $c=\cos\theta$, $\mu=M/m$. Separation when the wedge stops accelerating, $d(u^2)/dc=0$, giving $\mu=(2-3c)/c^3$. At $\theta=60^\circ$, $\mu=4$}, difficulty=10, diff_parts={K3 S3 I2 C2 F1}, diff_note={concepts: constraint, energy, separation condition | steps: constraint of load and wedge speeds; energy; maximise wedge speed; condition at separation; mass ratio | calculus on the energy}, group=TOPS04-P1-PHY-S3-G5-6]
For what mass ratio $\dfrac Mm$ will the rod form an angle $\theta=\pi/3$ with the vertical at the moment the load separates from the wedge?
\end{bankq}
\begin{bankq}[id=TOPS04-P1-PHY-S3-Q06, source={TOPS - Test 04 (18 Dec 2022)}, paper=1, qno=6, type=numerical, subtopic={Energy}, answer={1.414}, answer_src=key, difficulty=7, diff_parts={K2 S2 I2 C1 F1}, diff_note={concepts: constraint, energy | steps: ratio found; energy at the angle; wedge speed}, group=TOPS04-P1-PHY-S3-G5-6]
What is speed (in m/s) of the wedge at that moment?
\end{bankq}
\begin{bankq}[id=TOPS04-P2-PHY-S1-Q06, source={TOPS - Test 04 (18 Dec 2022)}, paper=2, qno=6, type=single, subtopic={Energy}, answer={C}, answer_src=key, difficulty=2, diff_parts={K1 S1 I0 C1 F0}, diff_note={concepts: momentum and kinetic energy | steps: energy in terms of momentum; percent change}]
If the momentum of a body is increased by 25\% the kinetic energy is increased by :-
\begin{qoptions}
\opt{A}{25\%}
\opt{B}{5\%}
\opt{C}{56\%}
\opt{D}{38\%}
\end{qoptions}
\end{bankq}
\begin{bankq}[id=TOPS04-P2-PHY-S1-Q07, source={TOPS - Test 04 (18 Dec 2022)}, paper=2, qno=7, type=single, subtopic={Potential energy curves}, answer={B}, answer_src=key, difficulty=5, diff_parts={K2 S2 I1 C1 F0}, diff_note={concepts: potential energy, energy conservation | steps: speed at the minimum of potential; energy; value}]
The potential energy of a 1 kg particle free to move along the $x$-axis is given by $V_x=\left(\dfrac{x^4}4-\dfrac{x^2}2\right)$ J. The total mechanical energy is 2J. Then the maximum speed in ms$^{-1}$ is :-
\begin{qoptions}
\opt{A}{2}
\opt{B}{$\dfrac3{\sqrt2}$}
\opt{C}{$\sqrt2$}
\opt{D}{$\dfrac1{\sqrt2}$}
\end{qoptions}
\end{bankq}
\begin{bankq}[id=TOPS04-P2-PHY-S1-Q08, source={TOPS - Test 04 (18 Dec 2022)}, paper=2, qno=8, type=single, subtopic={Work}, answer={C}, answer_src=key, difficulty=5, diff_parts={K2 S2 I1 C1 F0}, diff_note={concepts: variable force, integration | steps: velocity from the integral; first stop time; distance}]
A particle of mass $m$ is acted upon by a force $F=F_0\cos\left(\dfrac\pi2t\right)$. The distance travelled by the particle when it stops first (initial velocity of the particle is zero) is :
\begin{qoptions}
\opt{A}{$\dfrac{2F_0}{m\pi}$}
\opt{B}{$\dfrac{4F_0}{m\pi^2}$}
\opt{C}{$\dfrac{8F_0}{m\pi^2}$}
\opt{D}{zero}
\end{qoptions}
\end{bankq}
\begin{bankq}[id=TOPS04-P2-PHY-S1-Q09, source={TOPS - Test 04 (18 Dec 2022)}, paper=2, qno=9, type=single, subtopic={Work}, answer={D}, answer_src=key, difficulty=6, diff_parts={K2 S2 I2 C1 F0}, diff_note={concepts: work along a path, integration | steps: parametrise the path; component of the force; integrate; value}]
By applying a force $\vec F=(3xy-5z)\hat j+4z\hat k$ a particle is moved along the path $y=x^2$ from point $(0,0,0)$ to the point $(2,4,0)$. The work done by the $F$ on the particle is :-
\bfig{79.3}{TOPS04_p2_phy_s1_q09_1.png}
\begin{qoptions}
\opt{A}{$\dfrac{280}5$}
\opt{B}{$\dfrac{140}5$}
\opt{C}{$\dfrac{232}5$}
\opt{D}{$\dfrac{192}5$}
\end{qoptions}
\end{bankq}
\begin{bankq}[id=TOPS05-P1-PHY-S3-Q03, source={TOPS - Test 05 (14 Jan 2023)}, paper=1, qno=3, type=numerical, subtopic={Work}, answer={8.00}, answer_src=key, difficulty=6, diff_parts={K2 S2 I1 C1 F1}, diff_note={concepts: Atwood system, work by tension | steps: acceleration; tension; displacement in one second; work}]
A light inextensible string that goes over a smooth fixed pulley as shown in the figure connects two blocks of masses 0.36 kg and 0.72 kg. Taking $g=10$ m/s$^2$, find the work done (in joules) by the string on block of mass 0.36 kg during the first second after the system is released from rest. ($g=10$ m/s$^2$)
\bfig{47.5}{TOPS05_p1_phy_s3_q03_1.png}
\end{bankq}
\begin{bankq}[id=TOPS05-P2-PHY-S3-Q07, source={TOPS - Test 05 (14 Jan 2023)}, paper=2, qno=7, type=numerical, subtopic={Energy}, answer={1}, answer_src=key, difficulty=9, diff_parts={K2 S3 I2 C2 F1}, diff_note={concepts: energy conservation, string constraint | steps: geometry of the rise; constraint; energy balance; relation with height 4 m; mass ratio}]
In the figure, the ring is placed in the mid point and released. If $M$ rises to a maximum height of 4m. Find $M/m$. Assume no friction between string and ring.
\bfig{97.9}{TOPS05_p2_phy_s3_q07_1.png}
\end{bankq}

\section*{TMAS papers}
\begin{bankq}[id=TMAS02-P2-PHY-S1-Q05, source={TMAS - Test 02 (06 Aug 2023)}, paper=2, qno=5, type=multi, subtopic={Work}, answer={A,D}, answer_src=solved, working={Impulse on the vehicle is $\Delta p=5\times3=15$ N s; kinetic energy gained is 22.5 J, not 27 J; static friction accelerates the cart}, difficulty=7, diff_parts={K2 S2 I2 C1 F1}, diff_note={concepts: impulse, work and energy in a vehicle | steps: impulse; kinetic energy; potential energy change; friction role}]
For technology project, a student has built a vehicle, of total mass $5$ kg, that moves itself. As shown in the figure, it runs on four light wheels. A reel is attached to one of the axles, and a cord originally wound on the reel goes up over the pulley attached to the vehicle to support an elevated load. After the vehicle is released from rest, the load descends very slowly, unwinding the cord to turn the axle and make the vehicle move forward. Friction is negligible in pulley and axle bearings. The wheels do not slip on the floor. The reel has been constructed with a conical shape so that load descends at a constant low speed while the vehicle moves horizontally across the floor with a constant acceleration, reaching a final velocity of $3$ m/s. Mark the correct option(s) :
\bfig{45.1}{TMAS02_p2_phy_s1_q05_1.png}
\begin{qoptions}
\opt{A}{Magnitude of impulse on the vehicle due to friction of floor is $15$ N-s.}
\opt{B}{Work done by friction force between tyre and floor is $27$ J.}
\opt{C}{Change in gravitational potential energy of earth and cart system is $-27$ J.}
\opt{D}{The acceleration of the cart is due to friction force from the ground.}
\end{qoptions}
\end{bankq}
\end{document}
